\chapter{Fukaya 范畴}\label{chapter:II}
\SourcePage{103}
本章的主要目标，是在正合辛流形这一相对简单的特殊情形下，利用 Gromov 与 Floer 开创的经典分析方法，说明 Fukaya $\Ai$ 范畴的定义。可以说，这种方法如今已经过时，由 Fukaya–Oh–Ohta–Ono 的奠基性工作\bcite{60}、\bcite{56}所取代，后者允许在完全一般的情形下定义 Fukaya 范畴。不过，一般定义在分析与代数两方面都相当复杂；对于这里讨论的情形，其中许多精细技术并无必要。

为了方便非专门研究这一领域的读者，先从正合 Lagrangian 子流形的 Floer 上同调群的原始定义\bcite{44}开始。自 20 世纪 90 年代初以来，人们已知道 Floer 群上有一些运算，以适当一类具有边界标记点的 Riemann 曲面为指标，使之成为开弦拓扑量子场论（TQFT；这一想法似乎源于 Donaldson；讨论它或其闭弦对应物的一些论文有\bcite{119}、\bcite{106}、\bcite{135}、\bcite{131}）。在允许复结构变化的 Riemann 曲面族时，就出现了最初在\bcite{52}中提到的 Fukaya $\Ai$ 范畴（实际上，它只使用带边界标记点的圆盘，因此应视为某个更复杂理论的第一部分，而该理论仍有很大部分尚未探索）。这里所谓经典方法，是用零阶非齐次项扰动伪全纯映射方程，使解的模空间变得光滑，从而避免使用虚基本链。像\bcite{116}一样，非齐次项随 Riemann 曲面的复结构变化；我们将花一些篇幅说明精确要求。一旦澄清这一点，基本分析结果便没有特别困难或新颖之处，因此只给出尽量简短的证明概要。我们将更多强调分次与符号问题，因为两者对于应用都很重要。对于分次，采用 Kontsevich 的分次 Lagrangian 子流形形式体系\bcite{90}、\bcite{125}；与 Floer、Viterbo 较早提出的相对分次相比，它保留了稍多的信息。研究 Floer 理论中符号问题的一般框架建立于\bcite{48}，Lagrangian Floer 上同调的情形则在\bcite{60}中处理；但我们仍认为值得重新整理其中一部分理论。\znote{95}{两种分次的差别可以这样看：Floer--Viterbo 的相对分次只记录相位差落在 $\Z$ 里的部分，因此只能对每一对 $(L_0,L_1)$ 给出 $\Z$-分次，且需先任选一个基准；Kontsevich 的分次则给每个 $L$ 附加整体相位提升 $\theta_L:L\to\R$，于是分次是绝对的，平移 $[1]$ 在单个对象上也有定义。代价是这类提升并不总存在，通常要在 $2c_1(M)=0$ 之类的条件下工作；一旦存在，符号与次数的计算都可以逐点进行，这正是本节说它「保留了稍多信息」的意思。[来源：原书 §8 引言提到的 Kontsevich 分次 Lagrangian 形式体系；并见原书第三章关于分次与符号的讨论。]}

本章还讨论带对合的辛流形的等变 Fukaya 范畴，以及曲面的 Fukaya 范畴所具有的部分组合性质。后一主题已经在多种背景下出现：椭圆曲线的同调镜像对称\bcite{90}、\bcite{108}、\bcite{107}，曲面上曲线的 Floer 同调\bcite{135}、\bcite{35}，以及 Legendrian 纽结理论\bcite{30}、\bcite{43}。我们着重讨论将分析方法与组合方法联系起来的正则性论证，因为文献对此论述较少。
